% New two-dimensional local improvement. Existing files remain unchanged.
\subsection{A second-prefix upper-cap improvement in the fallback strip}
\label{sec:second-prefix-upper-fallback-20261009}

Set $\mu=25/128$, $\rho=19/96$, $\Lambda=153/256$ and
$v(t)=\min((1+t)/2,57/64-3t/2)$. The extended reference tree has
first stop $\rho$ and, for $\mu<x=t_1<\rho$, third endpoint exactly
$v(x)$; see the definition of $A$ in \texttt{extended-prefix-upper.tex}
and the common-reference paragraph of the joint density audit.
It does not use $v(t_2)$ in this strip.

Choose the witness-restricted realization of the already certified
higher actions: keep their fourth, fifth and sixth upgrades only on
their stated witness domains below $t_1=\mu$. Their certificates
require no gain from a hypothetical full fifth strategy above $\mu$.
This realization has the same certified density upper bound and the
same sixth-history guarantees on its witnesses. It avoids counting
any unmeasured higher-prefix gain twice. The following new action
changes only $\mu<t_1<\rho$ and is also disjoint from the new mixed
first-restoration strip $t_1>\rho$.

\paragraph{The new action and the complete gates.}
All rank-two nodes $p_1p_2$ under an admitted first prime already
exist in the old upper tree. On
\[
 \mathcal R=\{\mu<x<\rho,\quad a(x)<y<x\},\qquad
 a(x)=2v(x)-1=25/32-3x,
\]
replace their ordinary upper residual total level $v(x)$ by $v(y)$,
or by fixed strictly smaller constants in the finite construction.
This is precisely the region $v(y)>v(x)$: below $\mu$ the function
$v(y)$ is increasing, and above $\mu$ it is decreasing.

For $x\le\rho$ in this strip, $v(x)\ge3x$. Consequently on
$\mathcal R$,
\[
 2x< v(x),\qquad x+2y< v(x),\qquad
 2x+y< v(x)<v(y).
\]
The first two are the initial full-square conditions; the last is
the additional remaining-square gate when the selected factor is
$p_2$. Select $p_2$ and apply the selected-product factorization
lemma, including its capacity-below-one repair. Every later selected
remaining-square condition is exactly the corresponding full square.
The new ordinary upper odd cubes imply these squares and their next
ordered even squares. Thus this action requires neither a new
distribution theorem nor the two-large-prime lemma of the other note.

At rank two the node is retained, with coefficient $+1$. Its future
ordinary upper residual has odd first exclusions, hence full odd
first exclusions. The next rank-three cube may change, including
restoration of formerly excluded third nodes. Keeping the old
rank-three cube as an additional test would lose the claimed action.
Since $v(y)\le\Lambda<\Theta=299/500$, the new cube is compatible
with the original global upper support. Earlier admission decisions
are unchanged. This is a pointwise causal tree statement, not a
monotonicity assertion about finite signed density sums.

\paragraph{Density identity and its finite-tree scope.}
Write $\bar F=F/(2e^\gamma)$. The even-prefix identity gives the
limiting density saving, in the same $C_N/\log N$ units as $K$,
\begin{align*}
 D_2&=4\int_{\mathcal R}
 \frac{\bar F((v(x)-x-y)/y)-\bar F((v(y)-x-y)/y)}{xy^2}
 \,dx\,dy\\
 &=4\int_{\mathcal R}
 \frac{v(y)-v(x)}{xy(v(x)-x-y)(v(y)-x-y)}\,dx\,dy.
\end{align*}
Indeed both arguments are strictly between one and three. The lower
bound one follows from $v(x)>x+2y$; the upper follows already from
$v(y)\le\Lambda$, $x\ge\mu$, $y\ge3/16$, giving an upper bound
$55/48<3$. Hence the exact formula $\bar F(s)=1/s$ applies everywhere.
There is no ordered-slot factorial to insert.

The statement that the old rank-two residual is ordinary upper
support at total level $v(x)$ refers to the endpoint tree. In the
actual finite construction the third cap was disabled when the
third prime lay below the rough cutoff $N^\omega$. It is therefore
incorrect to claim the displayed identity for every arbitrary frozen
finite coarse tree without an error. The established extended-prefix
tail controls that difference by $O(\omega\log^2(1/\omega))$.
Use a common finite approximation, preserving the old witness
strategies below $\mu$, and let this error tend to zero in the stated
fixed-parameters-before-$N$ order. The same density identity then
holds in the endpoint limit. Every strictly smaller certified saving
is realizable with finite strict buffers and finite compilation error.

\paragraph{Two charts cover the whole improvement domain.}
Put $h=\rho-\mu=1/384$. Apart from null boundaries, the two disjoint
charts $k=-3$ and $k=1$ are
\[
 x=\mu+hu,\qquad y=\mu+khuv,\qquad 0<u,v<1.
\]
The negative chart covers $a(x)<y<\mu$ and the positive chart covers
$\mu<y<x$. In either chart
\[
 v(x)=\Lambda-\tfrac32hu,\qquad
 v(y)=\Lambda-\tfrac32huv,\qquad
 |\partial(x,y)/\partial(u,v)|=|k|h^2u.
\]
Define
\[
 R_o=\frac{53}{256}-hu(5/2+kv),\qquad
 R_n=\frac{53}{256}-hu(1+(k+3/2)v).
\]
Thus
\begin{equation}\label{eq:second-prefix-chart-integral}
 D_2=\sum_{k=-3,1}6|k|h^3
 \int_0^1\int_0^1\frac{u^2(1-v)}{xyR_oR_n}\,dv\,du.
\end{equation}
All four denominator factors are positive and bilinear or affine
on the closed square. There are no logarithms or uncomputed sieve
function values in this integral.

\paragraph{A closed rational certificate needing only two rectangles.}
In the negative chart,
$x\le\rho$, $y\le\mu$ and $R_o,R_n\le5/24$.
In the positive chart,
$x,y\le\rho$ and $R_o,R_n\le53/256$.
Since $\int_0^1\int_0^1u^2(1-v)\,dv\,du=1/6$,
\[
 D_{2,-}\ge\frac{3h^3}{\rho\mu(5/24)^2}=\frac3{95000},
 \qquad
 D_{2,+}\ge\frac{h^3}{\rho^2(53/256)^2}=\frac{32}{3042147}.
\]
Therefore, independently of any fine-grid output,
\begin{equation}\label{eq:second-prefix-coarse-positive}
 \boxed{D_2\ge\frac{640339}{15210735000}>4.2\cdot10^{-5}.}
\end{equation}

\paragraph{A bounded finer certificate, not an additional gain.}
The companion script subdivides each chart into $32^2$ rectangles.
For a chart rectangle $B=[u_0,u_1]\times[v_0,v_1]$, compute the maximum
of each denominator factor at its four corners. Bilinearity makes
these exact separate upper bounds. Let their product be $D_B^+$.
The exact positive numerator mass is
\[
 M_B=\frac{u_1^3-u_0^3}{3}
       \left(v_1-v_0-\frac{v_1^2-v_0^2}{2}\right).
\]
Store
\[
 b_B=\left\lfloor 2^{96}\,
                 \frac{6|k|h^3 M_B}{D_B^+}\right\rfloor.
\]
Then $2^{-96}\sum_B b_B$ is a rigorous lower bound for the same
$D_2$. The 2048 rows record the chart and indices, rational upper
denominator, exact numerator mass and integer floor. All arithmetic
is rational and floors are downward. Coarse and fine certificates
are alternative estimates of the same integral and must not be added.

\paragraph{Finite constants, boundaries and quantifiers.}
First trim the charts to $\delta\le u,v\le1-\delta$ with fixed
$\delta>0$. All branch improvements, selected gates, inherited
first-stop margins and residual arguments then have positive uniform
slack. Refine in the actual $(x,y)$ coordinates and choose, on each
retained small box, a constant below $\inf v(y)$ but above
$\sup v(x)$ and every required gate. Rank-two nodes are never deleted:
on unresolved boxes retain the old residual. The chart images can
be approximated from within by disjoint such boxes; their boundaries
are fixed affine curves. The positive rational kernel is bounded
on both closed charts, so the trimmed and boundary losses vanish.
The older finite first stop near $\rho$ and the transition near $\mu$
are handled within these same margins. Only afterwards fix the rough
cutoff and fine grid and take $N\to\infty$. The old tail, compiler and
positive-counting arguments apply to these actual signed coefficients.

Consequently the new endpoint upper bound is the previous certified
endpoint bound minus any chosen lower bound for $D_2$, with the usual
strict final allowance for one finite strategy. This note makes no
new global rounded $K$, $c$ or $b$ claim. No earlier source, state,
manuscript, formalization or publication artifact is modified.
